\documentclass[12pt, letterpaper]{article}

% Paquetes básicos
\usepackage[utf8]{inputenc}
\usepackage[T1]{fontenc}
\usepackage[spanish, es-tabla]{babel}
\usepackage[margin=2.5cm]{geometry}
\usepackage{hyperref}
\usepackage{enumitem}
\usepackage{titlesec}


% Formato de enlaces
\hypersetup{
    colorlinks=true,
    linkcolor=blue,
    filecolor=magenta,      
    urlcolor=cyan,
    pdftitle={Documentación IEEE 829 - SBAC},
}

\begin{document}

\begin{center}
    \vspace*{1cm}
    
    {\Large \textbf{Universidad Nacional Autónoma de México} \par}
    \vspace{0.5cm}
    {\large \textbf{Facultad de Ciencias} \par}
    
    \vspace{2cm}
    
    {\LARGE \textbf{Documentación de Pruebas IEEE 829} \par}
    \vspace{0.5cm}
    {\Large Proyecto Final: Sistema Básico de Administración de Configuración (SBAC) \par}
    
    \vspace{2cm}
    
    \textbf{Materia:} Pruebas de Software y Administración de la Configuración \\
    \textbf{Profesor:} Luis Rey Rutiaga Robles \\
    
    \vspace{1.5cm}
    
    \textbf{Integrantes:}\\
    Maryam Michelle Del Monte Ortega \\
    José Luis Flores Gutiérrez \\
    Erik Eduardo Gómez López\\
    Carlos Adrián Celaya Nava \\

    \vspace{2cm}
    
    \textbf{Fecha:}\\
    04 de junio de 2026
    
    \vfill
\end{center}

\newpage

\section{Plan de Pruebas}
\begin{itemize}[leftmargin=*]
    \item \textbf{Identificador:} TP-SBAC-1.0
    \item \textbf{Introducción:} Este documento explica de forma clara cómo pusimos a prueba nuestro Sistema Básico de Administración de Configuración (SBAC). Nuestro objetivo es asegurar que esta herramienta de terminal (CLI) funcione sin problemas al guardar las versiones de los archivos.
    \item \textbf{?`Qué vamos a probar?:}
    \begin{itemize}
        \item Que el repositorio se cree y se lea correctamente (\texttt{init}, \texttt{status}).
        \item Que el sistema detecte y guarde los archivos sin perder información (\texttt{add}, \texttt{commit}).
        \item Que el historial y las etiquetas sean fáciles de consultar (\texttt{history}, \texttt{baseline}, \texttt{list-baselines}, \texttt{diff}).
        \item Que podamos viajar en el tiempo para recuperar archivos viejos (\texttt{checkout}).
    \end{itemize}
    \item \textbf{?`Cómo lo vamos a probar?:} Usaremos un enfoque mixto. Revisaremos que el código interno haga lo que debe (caja blanca) y probaremos el programa como si fuéramos un usuario final (caja negra). Todo esto está automatizado con las herramientas \texttt{unittest} y \texttt{pytest}, y se ejecuta dentro de un contenedor Docker (\texttt{python:3.11-slim}) para que funcione igual en cualquier computadora.
    \item \textbf{Criterios de Éxito:} Consideramos que el sistema está listo si logra pasar automáticamente las 15 pruebas diseñadas y si la herramienta \texttt{pytest-cov} nos confirma que estamos evaluando más del 95\% de todo nuestro código escrito.
\end{itemize}

\section{Especificaciones de Diseño de Pruebas}
\begin{itemize}[leftmargin=*]
    \item \textbf{Identificador:} TDS-SBAC-1.0
    \item \textbf{Puntos clave a revisar:} Queremos estar seguros de que la carpeta \texttt{.sbac} no se corrompa, que el archivo de configuración \texttt{metadata.json} guarde bien los datos, y que el programa no colapse si el usuario escribe un comando equivocado.
    \item \textbf{Nuestra forma de trabajo:}
    \begin{itemize}
        \item \textbf{División de tareas:} Separamos las pruebas en tres grupos: pruebas pequeñas para funciones individuales (unitarias), pruebas para ver cómo se comunican las funciones entre sí (integración), y pruebas para asegurar que los archivos físicos viajen bien en el tiempo (regresión).
        \item \textbf{Entorno seguro:} Cada vez que corremos una prueba, creamos una carpeta temporal invisible (\texttt{tempfile.mkdtemp()}). Así podemos probar cómo se copian y borran los archivos sin poner en riesgo nuestro proyecto real.
    \end{itemize}
\end{itemize}

\section{Casos de Prueba}
Aquí describimos los escenarios más importantes que evaluamos. \textit{(Nota: Nuestro sistema automatizado corre un total de 15 pruebas para cubrir estos casos y varios más)}.

\begin{itemize}[leftmargin=*]
    \item \textbf{TC-01: Iniciar el repositorio por primera vez}
    \begin{itemize}
        \item \textbf{?`Qué necesitamos antes?:} Nada, solo una carpeta vacía.
        \item \textbf{Pasos:} Escribir el comando \texttt{sbac start} o \texttt{sbac init}.
        \item \textbf{Qué debería pasar:} El sistema debe crear una carpeta oculta \texttt{.sbac/}, un espacio para guardar los archivos (\texttt{snapshots/}) y un documento JSON listo para registrar los cambios.
    \end{itemize}
    
    \item \textbf{TC-02: Evitar que el programa colapse por descuidos del usuario}
    \begin{itemize}
        \item \textbf{?`Qué necesitamos antes?:} Estar en una carpeta donde NO se haya iniciado SBAC.
        \item \textbf{Pasos:} Intentar usar comandos como \texttt{sbac add} o \texttt{sbac commit}.
        \item \textbf{Qué debería pasar:} El programa debe detenerse amablemente y avisarle al usuario que primero debe iniciar el repositorio, en lugar de mostrar un error feo de Python en la pantalla.
    \end{itemize}
    
    \item \textbf{TC-03: Comparar versiones usando etiquetas (Baselines)}
    \begin{itemize}
        \item \textbf{?`Qué necesitamos antes?:} Tener un archivo modificado y guardado con dos etiquetas distintas (por ejemplo, ``V1'' y ``V2'').
        \item \textbf{Pasos:} Escribir \texttt{sbac diff V1 V2}.
        \item \textbf{Qué debería pasar:} El sistema debe entender qué versión significa cada etiqueta y mostrarnos en pantalla exactamente qué líneas de código se agregaron o borraron.
    \end{itemize}
    
    \item \textbf{TC-04: Viajar al pasado (Checkout)}
    \begin{itemize}
        \item \textbf{?`Qué necesitamos antes?:} Tener un archivo guardado, haberlo modificado, y haber guardado esa nueva versión.
        \item \textbf{Pasos:} Escribir \texttt{sbac checkout v1} para pedirle que regrese al estado original.
        \item \textbf{Qué debería pasar:} El archivo actual debe sobrescribirse para volver a ser idéntico a la versión antigua que le pedimos.
    \end{itemize}
\end{itemize}

\section{Procedimientos de Prueba}
\begin{itemize}[leftmargin=*]
    \item \textbf{Identificador:} TPR-SBAC-1.0
    \item \textbf{Requisitos Previos:} Solo necesitas tener Docker instalado en tu computadora.
    \item \textbf{Pasos para correr las pruebas:}
    \begin{enumerate}
        \item Abre tu terminal en la carpeta del proyecto SBAC.
        \item Empaqueta el entorno de pruebas escribiendo: \\ \texttt{docker build -t sbac-tests .}
        \item Inicia las pruebas automatizadas escribiendo: \\ \texttt{docker run --rm sbac-tests}
        \item Observa la pantalla: verás la confirmación de que todas las pruebas pasaron en color verde y una tabla mostrando el porcentaje de código evaluado.
    \end{enumerate}
\end{itemize}

\section{Informes de Incidentes}
\begin{itemize}[leftmargin=*]
    \item \textbf{Identificador:} TIR-SBAC-01
    \item \textbf{Fecha del problema:} 29 de mayo de 2026.
    \item \textbf{?`Qué pasó?:} Durante las primeras pruebas, notamos que si intentábamos comparar dos versiones usando sus nombres de etiqueta (como \texttt{ESTABLE}), el programa colapsaba y se cerraba. Solo funcionaba si usábamos los números internos de versión (como \texttt{v2}).
    \item \textbf{Problema causado:} No estábamos cumpliendo con una de las funciones principales que prometía el manual del usuario para visualizar cambios entre líneas base.
    \item \textbf{Cómo lo arreglamos:} Modificamos la función \texttt{show\_diff}. Ahora, antes de comparar, el sistema revisa si la palabra que escribió el usuario es una etiqueta, si lo es, la traduce silenciosamente a su número de versión interno y la comparación funciona perfecto. \textbf{Incidente cerrado.}
\end{itemize}

\vspace{0.3cm}

\begin{itemize}[leftmargin=*]
    \item \textbf{Identificador:} TIR-SBAC-02
    \item \textbf{Fecha del problema:} 30 de mayo de 2026.
    \item \textbf{?`Qué pasó?:} Al realizar pruebas de estrés en la función de seguimiento, detectamos que si un usuario ejecutaba \texttt{sbac add archivo.txt} múltiples veces seguidas, el archivo se registraba repetidamente en el arreglo del \texttt{metadata.json}.
    \item \textbf{Problema causado:} Duplicidad de datos. Esto provocaba que al hacer un \texttt{commit}, el sistema intentara copiar el mismo archivo físico varias veces, consumiendo recursos innecesarios y causando comportamientos erráticos al leer el historial.
    \item \textbf{Cómo lo arreglamos:} Se implementó una validación condicional en la función \texttt{add\_file} (\texttt{if filename not in metadata[``tracked\_files'']:}). Ahora el sistema verifica si el archivo ya está en seguimiento y, de ser así, omite el registro y avisa al usuario amablemente. \textbf{Incidente cerrado.}
\end{itemize}

\vspace{0.3cm}

\begin{itemize}[leftmargin=*]
    \item \textbf{Identificador:} TIR-SBAC-03
    \item \textbf{Fecha del problema:} 1 de junio de 2026.
    \item \textbf{?`Qué pasó?:} Durante las pruebas de caminos alternos (edge cases), se ejecutó el comando \texttt{sbac commit ``Mensaje''} en un repositorio recién inicializado que no tenía ningún archivo agregado. 
    \item \textbf{Problema causado:} Generación de ``commits fantasma''. El sistema creaba un nuevo ID de versión y generaba una carpeta vacía dentro de \texttt{snapshots/}, alterando el orden del historial sin guardar información real.
    \item \textbf{Cómo lo arreglamos:} Se agregó una barrera preventiva al inicio de la función \texttt{create\_commit}. Si la lista de archivos en seguimiento está vacía, la ejecución se aborta inmediatamente y se notifica al usuario que debe usar \texttt{sbac add} primero. \textbf{Incidente cerrado.}
\end{itemize}

\section{Informes de Prueba}
\begin{itemize}[leftmargin=*]
    \item \textbf{Identificador:} TSR-SBAC-1.0
    \item \textbf{Resumen del trabajo:} Corrimos todo el paquete de pruebas sobre la versión final del sistema sin ningún contratiempo. Para tener el código más limpio y ordenado, separamos las pruebas en tres archivos distintos (\texttt{test\_unit.py}, \texttt{test\_integration.py} y \texttt{test\_regression.py}).
    \item \textbf{Nuestros Resultados:}
    \begin{itemize}
        \item \textbf{Pruebas realizadas:} 15
        \item \textbf{Pruebas que pasaron exitosamente:} 15 (100\% de éxito).
        \item \textbf{Porcentaje de código evaluado:} 99\% (150 de 152 líneas de código).
    \end{itemize}
    \item \textbf{?`Por qué sacamos 99\% y no 100\%?:} Alcanzar un 99\% es un indicador de calidad excelente. El 1\% faltante son solo dos líneas de código interno (\texttt{132} y \texttt{134}) que actúan como ``seguros de vida'' al momento de traducir las etiquetas de las versiones. Es muy difícil obligar al entorno de pruebas a activar estas alarmas preventivas, lo cual es algo completamente normal y aceptado en la industria del software.
    \item \textbf{Conclusión Final:} El sistema SBAC cumple con su propósito de guardar versiones de forma confiable. Demostró ser muy resistente incluso cuando intentamos engañarlo borrando archivos a la fuerza o pidiéndole datos vacíos. Al usar Docker, garantizamos que el profesor o cualquier usuario pueda usarlo sin tener que configurar nada en su computadora. El software está listo para entregarse.
\end{itemize}

\end{document}
